\section{Implementation Consistency Lemmas}
\label{sec:consistency-lemmas}

These lemmas verify that the recursive implementation and its accumulated
representation agree internally. They do not introduce new mathematical
properties but are essential for formal software consistency.

\begin{itemize}
  \item Element consistency: $I_k = acc_k$ --- Section~5.2
  \item Accumulated delta consistency: $acc_{p+1} - acc_p = L_{p+1}$
    --- Section~5.3
  \item Last element agreement: $\operatorname{last}(I) = acc_{n-1} =
    I_{n-1}$ --- Section~5.4
  \item Size agreement: $\lvert acc\rvert = \lvert L\rvert$ ---
    Section~5.5
\end{itemize}

\subsection{Accumulated List Definition}
\label{sec:accumulated-list-definition}

The accumulated list represents the discrete integral as a full list of
partial sums rather than element-by-element access.

Let:

\begin{equation*}
\begin{aligned}
& \operatorname{acc}(L, init) \in \mathbb{Z}^{\lvert L\rvert} \\
& L = [x_0, x_1, \dots, x_{n-1}]
\end{aligned}
\end{equation*}

Then, the accumulated list is defined recursively as:

\begin{equation*}
\operatorname{acc}(L, init) :=
\begin{cases}
L_e & \text{if } L = L_e \\
(\operatorname{head}(L) + init) :: \operatorname{acc}(\operatorname{tail}(L),\ \operatorname{head}(L) + init) & \text{otherwise}
\end{cases}
\end{equation*}

The full Integral implementation including the \texttt{acc} method is at
\href{https://github.com/thiagomata/prime-numbers/blob/master/src/main/scala/v1/chapter3/list/integral/Integral.scala}{\texttt{Integral.scala}}:

\begin{lstlisting}[style=scala]
case class Integral(list: List[BigInt], init: BigInt = 0) {
  def apply(position: BigInt): BigInt = {
    require(list.nonEmpty)
    require(position >= 0 && position < list.size)
    if (position == 0) this.head else Integral(list.tail, this.head).apply(position - 1)
  }
  def acc: List[BigInt] = {
    decreases(list.size)
    if (list.isEmpty) list else List(this.head) ++ Integral(list.tail, this.head).acc
  }
  def head: BigInt = {
    require(list.nonEmpty)
    list.head + init
  }
  def tail: List[BigInt] = {
    require(list.nonEmpty)
    Integral(list.tail, this.head).acc
  }
  def last: BigInt = {
    require(list.nonEmpty)
    acc.last
  }
}
\end{lstlisting}

\subsection{Element Consistency}
\label{sec:element-consistency}

The $k\text{-th}$ element of the Integral equals the $k\text{-th}$ element
of the accumulated list.

\begin{equation*}
\forall\, k \in [0, n-1]:\ I_k = acc_k
\end{equation*}

\begin{equation*}
\begin{aligned}
& L &= x_0 :: \operatorname{tail}(L)
  & \qquad &\text{[List decomposition]} \\
& I &= (x_0 + i) :: \operatorname{tail}(I)
  & \qquad &\text{[Integral decomposition]} \\
& \operatorname{acc}(L, i) &= (x_0 + i) :: \operatorname{acc}(\operatorname{tail}(L),(x_0 + i))
  & \qquad &\text{[Definition of } \operatorname{acc}] \\
& I_0 &= x_0 + i = \operatorname{acc}_0
  & \qquad &\text{[Base case]} \\
& I_{(p+1)} &= \operatorname{tail}(I)_p
  & \qquad &\text{[Tail Access Shift Left]} \\
& \operatorname{acc}_{(p+1)} &= \operatorname{acc}(\operatorname{tail}(L),(x_0 + i))_p
  & \qquad &\text{[Recursive accumulation]} \\
& \operatorname{tail}(L)_p &= \operatorname{acc}(\operatorname{tail}(L), (x_0 + i))_p
  & \qquad &\text{[Inductive hypothesis]} \\
& \implies \quad I_{p+1} &= \operatorname{acc}_{p+1}
  & \qquad &\text{[By substitution]} \\
&& \therefore \\
& \forall p \in [0..n-1], \quad I_p &= \operatorname{acc}_p \quad \blacksquare
  & \qquad &\text{[Q.E.D.]}
\end{aligned}
\end{equation*}

{\raggedright
This property is verified in the
\href{https://github.com/thiagomata/prime-numbers/blob/master/src/main/scala/v1/chapter3/list/integral/properties/IntegralProperties.scala}{\texttt{IntegralProperties::\allowbreak assertAccMatchesApply}}.
The full Scala verification code is in Appendix~A.5.
\par}

\subsection{Accumulated Delta Consistency}
\label{sec:accumulated-delta-consistency}

The difference between two consecutive accumulated values in Acc equals
the corresponding value from the original list.

\begin{equation*}
\forall\ p \in [0, n-2]:\ \operatorname{acc}_{p+1} - \operatorname{acc}_p = L_{p+1}
\end{equation*}

\begin{equation*}
\begin{aligned}
& L &= [x_0, x_1, \dots, x_{n-1} ]
  & \qquad &\text{[List definition]} \\
& L &= x_0 :: \operatorname{tail}(L)
  & \qquad &\text{[List decomposition]} \\
& \operatorname{acc}(L, i) &= (x_0 + i) :: \operatorname{acc}(\operatorname{tail}(L), x_0 + i)
  & \qquad &\text{[Definition of acc]} \\
& \operatorname{acc}_0 &= x_0 + i
  & \qquad &\text{[Base case]} \\
& \operatorname{acc}_1 &= x_1 + \operatorname{acc}_0
  & \qquad &\text{[Recursive accumulation]} \\
& \implies \quad \operatorname{acc}_1 - \operatorname{acc}_0 &= x_1 = L_1
  & \qquad &\text{[Cancellation]} \\
& \operatorname{acc}_{p+1} &= x_{p+1} + \operatorname{acc}_p
  & \qquad &\text{[Recursive accumulation]} \\
& \implies \quad \operatorname{acc}_{p+1} - \operatorname{acc}_p &= x_{p+1} = L_{p+1}
  & \qquad &\text{[By subtraction]}
\end{aligned}
\end{equation*}

\begin{equation*}
\therefore
\end{equation*}

\begin{equation*}
\forall p \in [0..n-2],\quad \operatorname{acc}_{p+1} - \operatorname{acc}_p
  = L_{p+1} \quad \blacksquare \qquad \text{[Q.E.D.]}
\end{equation*}

{\raggedright
This property is verified in the
\href{https://github.com/thiagomata/prime-numbers/blob/master/src/main/scala/v1/chapter3/list/integral/properties/IntegralProperties.scala}{\texttt{IntegralProperties::\allowbreak assertAccDiffMatchesList}}.
The full Scala verification code is in Appendix~A.6.
\par}

\subsection{Last Element Agreement}
\label{sec:last-element-agreement}

The last element of the accumulated list equals the last element of the
integral, which is the element at position $n-1$.

\begin{equation*}
\begin{aligned}
acc_{(n - 1)} & = \operatorname{last}(I) \\
acc_{(n - 1)} & = I_{(n - 1)}
\end{aligned}
\end{equation*}

\begin{equation*}
\begin{aligned}
& L &= [x_0, x_1, \dots, x_{n-1}]
  & \qquad &\text{[List definition]} \\
& \operatorname{last}(L) &=
  \begin{cases}
  x_0 & \text{if } \lvert L\rvert = 1 \\
  \operatorname{last}(\operatorname{tail}(L)) & \text{if } \lvert L\rvert > 1
  \end{cases}
  & \qquad &\text{[Definition of last]} \\
& \operatorname{acc}(L, i) &= (x_0 + i) :: \operatorname{acc}(\operatorname{tail}(L), x_0 + i)
  & \qquad &\text{[Definition of accumulation]} \\
& I &= \operatorname{acc}(L, i)
  & \qquad &\text{[Integral as accumulated list]}
\end{aligned}
\end{equation*}

\subsubsection*{Base case: $\lvert L\rvert = 1$}

\begin{equation*}
\begin{aligned}
& L &= [x_0]
  & \qquad &\text{[Singleton list]} \\
& \operatorname{acc}(L, i) &= [x_0 + i]
  & \qquad &\text{[By definition]} \\
& I &= [x_0 + i]
  & \qquad &\text{[Integral is acc]} \\
& \operatorname{last}(I) &= x_0 + i = acc_0 = I_0
  & \qquad &\text{[last on singleton]}
\end{aligned}
\end{equation*}

\subsubsection*{Inductive step: $\lvert L\rvert > 1$}

\begin{equation*}
\begin{aligned}
& L &= x_0 :: \operatorname{tail}(L)
  & \qquad &\text{[List decomposition]} \\
& I &= (x_0 + i) :: \operatorname{acc}(\operatorname{tail}(L), x_0 + i)
  & \qquad &\text{[Recursive definition]} \\
& \operatorname{tail}(I) &= \operatorname{acc}(\operatorname{tail}(L), x_0 + i)
  & \qquad &\text{[Tail of integral]} \\
& \operatorname{last}(I) &= \operatorname{last}(\operatorname{tail}(I))
  & \qquad &\text{[Recursive last]} \\
& \operatorname{last}(\operatorname{tail}(I)) &= \operatorname{acc}(\operatorname{tail}(L), x_0 + i)_{(n - 2)}
  & \qquad &\text{[Inductive hypothesis]} \\
& &= acc_{(n - 1)}
  & \qquad &\text{[Shifted indexing]} \\
& \implies \ \operatorname{last}(I) &= acc_{(n - 1)} = I_{(n - 1)}
  & \qquad &\text{[By substitution]}
\end{aligned}
\end{equation*}

\begin{equation*}
\therefore
\end{equation*}

\begin{equation*}
\operatorname{last}(I) = acc_{(n - 1)} = I_{(n - 1)} \quad \blacksquare
  \qquad \text{[Q.E.D.]}
\end{equation*}

{\raggedright
This property is verified in the
\href{https://github.com/thiagomata/prime-numbers/blob/master/src/main/scala/v1/chapter3/list/integral/properties/IntegralProperties.scala}{\texttt{IntegralProperties::\allowbreak assertLast}}.
The full Scala verification code is in Appendix~A.7.
\par}

\subsection{Size Agreement}
\label{sec:size-agreement}

The size of the accumulated list equals the size of the original list.

\begin{equation*}
\lvert acc\rvert = \lvert L\rvert
\end{equation*}

\begin{equation*}
\begin{aligned}
& L &= [x_0, x_1, \dots, x_{n-1}]
  & \qquad &\text{[List definition]} \\
& \operatorname{acc}(L, i) &= (x_0 + i) :: \operatorname{acc}(\operatorname{tail}(L), x_0 + i)
  & \qquad &\text{[Recursive accumulation]}
\end{aligned}
\end{equation*}

\subsubsection*{Empty List: $\lvert L\rvert = 0$}

\begin{equation*}
\begin{aligned}
& L &= []
  & \qquad &\text{[Empty list]} \\
& \operatorname{acc}(L, i) &= []
  & \qquad &\text{[By definition]} \\
& \lvert \operatorname{acc}(L, i)\rvert &= 0 = \lvert L\rvert
  & \qquad &\text{[Equal size]}
\end{aligned}
\end{equation*}

\subsubsection*{Singleton List: $\lvert L\rvert = 1$}

\begin{equation*}
\begin{aligned}
& L &= [x_0]
  & \qquad &\text{[Singleton list]} \\
& \operatorname{acc}(L, i) &= [x_0 + i]
  & \qquad &\text{[By definition]} \\
& \lvert \operatorname{acc}(L, i)\rvert &= 1 = \lvert L\rvert
  & \qquad &\text{[Equal size]}
\end{aligned}
\end{equation*}

\subsubsection*{Inductive step: $\lvert L\rvert > 1$}

\begin{equation*}
\begin{aligned}
& L &= x_0 :: \operatorname{tail}(L)
  & \qquad &\text{[Decomposition]} \\
& \operatorname{acc}(L, i) &= (x_0 + i) :: \operatorname{acc}(\operatorname{tail}(L), x_0 + i)
  & \qquad &\text{[Recursive call]} \\
& \lvert \operatorname{acc}(\operatorname{tail}(L), x_0 + i)\rvert &= \lvert \operatorname{tail}(L)\rvert
  & \qquad &\text{[Inductive hypothesis]} \\
& \lvert \operatorname{acc}(L, i)\rvert &= 1 + \lvert \operatorname{tail}(L)\rvert = \lvert L\rvert
  & \qquad &\text{[Cons adds 1]} \\
& & \therefore \\
& \lvert \operatorname{acc}(L, i)\rvert &= \lvert L\rvert \quad \blacksquare
  & \qquad &\text{[Q.E.D.]}
\end{aligned}
\end{equation*}

{\raggedright
This property is verified in the
\href{https://github.com/thiagomata/prime-numbers/blob/master/src/main/scala/v1/chapter3/list/integral/properties/IntegralProperties.scala}{\texttt{IntegralProperties::\allowbreak assertSizeAccEqualsSizeList}}.
The full Scala verification code is in Appendix~A.8.
\par}
